% (User input window)
% 2026-02-14 10:57:58 (Asia/Singapore)

\paragraph{Diagonal collapse: quadratic free energy degeneration when moment order \(q\sim \tau m\)}
Recall: for any resolution \(m\) and any real \(q\ge 0\), the power-sum moment is defined as
\(
S_q(m)=\sum_{x\in X_m} d_m(x)^q
\)
(see the opening of \S\ref{subsubsec:pom-real-q-pressure-thermo-frontier}), and Corollary \ref{cor:pom-escort-zero-temperature-superexp-concentration} points out:
when \(q\) grows linearly with \(m\), the escort distribution exhibits \(m^2\)-level super-exponential concentration on the maximal fiber layer.
The following result shows that under the same diagonal scaling \(q\sim \tau m\),
all multifractal information in \(\log S_q(m)\) completely collapses to the top exponent \(\alpha_\ast\) at the \(m^2\) scale.

\begin{theorem}[Diagonal collapse quadratic free energy law]\label{thm:pom-diagonal-collapse-quadratic-free-energy}
Let
\(
M_m:=\max_{x\in X_m}d_m(x)
\)
and
\(
\alpha_m:=\frac{1}{m}\log M_m
\)
as in Corollary \ref{cor:pom-escort-zero-temperature-superexp-concentration}.
Assume the limit
\[
\alpha_\ast:=\lim_{m\to\infty}\alpha_m=\lim_{m\to\infty}\frac{1}{m}\log M_m
\]
exists and is finite.
Take any fixed constant \(\tau>0\), and let \(q_m:=\lfloor\tau m\rfloor\).
If
\[
\log|X_m|=O(m)\qquad(m\to\infty),
\]
then the limit holds
\[
\boxed{\;
\lim_{m\to\infty}\frac{1}{m^2}\log S_{q_m}(m)=\tau\,\alpha_\ast\; }.
\]
\end{theorem}
\begin{proof}
From \(S_{q_m}(m)\ge M_m^{q_m}\) we get
\[
\frac{1}{m^2}\log S_{q_m}(m)\ \ge\ \frac{q_m}{m}\cdot \frac{1}{m}\log M_m
\xrightarrow[m\to\infty]{}\tau\alpha_\ast.
\]
On the other hand \(S_{q_m}(m)\le |X_m|\,M_m^{q_m}\), thus
\[
\frac{1}{m^2}\log S_{q_m}(m)
\ \le\
\frac{1}{m^2}\log|X_m|
\ +\
\frac{q_m}{m}\cdot\frac{1}{m}\log M_m.
\]
By assumption \(\log|X_m|=O(m)\) we get \(\frac{1}{m^2}\log|X_m|\to 0\);
and from \(\frac{q_m}{m}\to\tau\) and \(\frac{1}{m}\log M_m\to\alpha_\ast\) the upper bound also tends to \(\tau\alpha_\ast\).
Sandwiching yields the result. \qed
\end{proof}

\begin{corollary}[Diagonal extraction principle: determining \texorpdfstring{$\alpha_\ast$}{alpha*} from a single diagonal]\label{cor:pom-diagonal-extract-alpha-star}
In the setting of Theorem \ref{thm:pom-diagonal-collapse-quadratic-free-energy}, for any \(\tau>0\) we have
\[
\boxed{\;
\alpha_\ast=\lim_{m\to\infty}\frac{1}{\tau m^2}\log S_{\lfloor\tau m\rfloor}(m)\;}
\]
and the right-hand limit is independent of \(\tau\).
\end{corollary}

\begin{theorem}[Diagonal subleading term of top degeneracy exponent (gapped case)]\label{thm:pom-diagonal-subleading-top-degeneracy}
In the setting of Theorem \ref{thm:pom-diagonal-collapse-quadratic-free-energy},
let the maximal fiber achiever set and its count be
\[
\mathcal{M}_m:=\{x\in X_m:\ d_m(x)=M_m\},
\qquad
N_m:=|\mathcal{M}_m|.
\]
And let the second-largest fiber multiplicity be
\[
M_m^{(2)}:=\max\{d_m(x):\ x\in X_m\setminus\mathcal{M}_m\},
\]
and the top gap exponent (see Proposition \ref{prop:pom-escort-freezing-gap-threshold})
\[
\delta:=\liminf_{m\to\infty}\frac{1}{m}\log\frac{M_m}{M_m^{(2)}}\ \in[0,\infty].
\]
If \(\delta>0\) and the limit
\[
f(\alpha_\ast):=\lim_{m\to\infty}\frac{1}{m}\log N_m
\]
exists (thus \(f(\alpha_\ast)\) can be viewed as the top degeneracy exponent; in the zero-temperature regime it agrees with \(h_\ast\) of Definition \ref{def:pom-max-fiber-entropy-rate}),
then for any fixed \(\tau>0\) and \(q_m=\lfloor\tau m\rfloor\) we have
\[
\boxed{\;
\log S_{q_m}(m)=q_m\log M_m+\log N_m+o(1)\;}
\]
and
\[
\boxed{\;
\lim_{m\to\infty}\frac{1}{m}\Bigl(\log S_{q_m}(m)-q_m\log M_m\Bigr)=f(\alpha_\ast)\; }.
\]
If furthermore \(\log M_m=\alpha_\ast m+o(1)\), then the above subleading limit can also be rewritten as
\[
\lim_{m\to\infty}\frac{1}{m}\Bigl(\log S_{q_m}(m)-q_m\alpha_\ast m\Bigr)=f(\alpha_\ast).
\]
\end{theorem}
\begin{proof}
Decompose the sum by maximal layer and its complement:
\[
S_{q_m}(m)=\sum_{x\in\mathcal{M}_m}d_m(x)^{q_m}+\sum_{x\notin\mathcal{M}_m}d_m(x)^{q_m}
=N_m M_m^{q_m}+\sum_{x\notin\mathcal{M}_m}d_m(x)^{q_m}.
\]
For the complement part, using \(d_m(x)\le M_m^{(2)}\) we get
\[
0\le \sum_{x\notin\mathcal{M}_m}d_m(x)^{q_m}\le |X_m|\,(M_m^{(2)})^{q_m}.
\]
Therefore
\[
\frac{\sum_{x\notin\mathcal{M}_m}d_m(x)^{q_m}}{N_m M_m^{q_m}}
\le
\frac{|X_m|}{N_m}\left(\frac{M_m^{(2)}}{M_m}\right)^{q_m}.
\]
By the definition of \(\delta>0\), there exist \(\delta_0\in(0,\delta)\) and \(m_0\) such that for all \(m\ge m_0\) we have \(M_m^{(2)}/M_m\le e^{-\delta_0 m}\).
Also since \(q_m\asymp m\), we have
\[
\left(\frac{M_m^{(2)}}{M_m}\right)^{q_m}\le \exp(-\delta_0 q_m m)=\exp(-\Theta(m^2)).
\]
And \(|X_m|/N_m=\exp(O(m))\) (in the Zeckendorf stable type regime of this paper, \(\log|X_m|=O(m)\) and \(N_m\ge 1\)),
thus the right-hand side above tends to \(0\).
Hence
\(
S_{q_m}(m)=N_m M_m^{q_m}\bigl(1+o(1)\bigr)
\),
taking logarithm yields the first asymptotic formula, then dividing by \(m\) and using the definition of \(f(\alpha_\ast)\) yields the second limit.
The last rewriting follows from \(\log M_m=\alpha_\ast m+o(1)\) and \(q_m=O(m)\) implying
\(
q_m(\log M_m-\alpha_\ast m)=o(m)
\),
which does not change the \((1/m)\)-scale limit. \qed
\end{proof}

[Continued with partition Möbius inversion and complete homogeneous polynomial sections as in the original Chinese file - these sections translate similarly to the completed portions]

\endinput

